\chapter{Durchführung} \label{ch:Durchfuehrung}

\section{Messverfahren} \label{sec:Messverfahren}

Vor Beginn jeder Messreihe wurde der Kreisel so justiert, dass er als kräftefrei gelten konnte.  
Dazu wurde die Farbscheibe auf den Aluminiumstab gesteckt und der Kreisel mit geringer Drehzahl in horizontale Lage gebracht.  
Durch Verschieben der Scheibe entlang des Stabs wurde die Lage so angepasst, dass keine Präzession mehr auftrat.

\subsection{Messung der Dämpfung}

Zur Bestimmung der Dämpfungskonstante wurde der Kreisel mit beiden Zusatzgewichten am Stabende auf etwa \SI{600}{\per\minute} bis \SI{700}{\per\minute} beschleunigt.  
Über einen Zeitraum von zwölf Minuten wurde alle zwei Minuten die Eigenfrequenz mithilfe des Tachometers bestimmt.  
Die logarithmische Darstellung der Kreisfrequenz erlaubt anschließend die Bestimmung der Dämpfungskonstante über einen linearen Fit.

\subsection{Messung der Präzession}

Für die Untersuchung des schweren Kreisels wurden Zusatzgewichte in verschiedenen Abständen (15 cm und 20 cm) angebracht.  
Der Kreisel wurde jeweils bei senkrechter Achse auf eine definierte Drehfrequenz gebracht und anschließend unter einem Winkel von \SI{45}{\degree} oder \SI{90}{\degree} ausgelenkt.  
Nach dem Loslassen wurde die Präzessionsdauer mit einer Stoppuhr bestimmt.  
Für jede Gewichtskonfiguration wurden drei Messungen bei unterschiedlichen Eigenfrequenzen durchgeführt.

\subsection{Messung der Nutation}

Zur Beobachtung der Nutation wurde der Kreisel kräftefrei betrieben und durch einen leichten Stoß am Kugellager in Nutation versetzt.  
Die momentane Drehachse wurde mithilfe der Farbscheibe sichtbar gemacht, indem der Farbwechsel im Zentrum der Scheibe beobachtet wurde.  
Für fünf verschiedene Eigenfrequenzen wurde die Zeit für zehn vollständige Farbwechsel gemessen, woraus sich die Umlauffrequenz $\Omega$ ergibt.

\img[.5]{momentare_drehachse}{Schematische Darstellung der farbigen Drehscheibe. \cite{skriptPAP21}}

Zusätzlich wurde die Nutationsfrequenz direkt mit dem Stroboskop bestimmt.  
Dazu wurde zunächst die Nutationsfrequenz eingestellt, bei der die Stange scheinbar stillsteht, und unmittelbar danach die Eigenfrequenz abgelesen.  
Da der Kreisel gedämpft ist, mussten beide Messungen schnell hintereinander erfolgen.
